\section{Conclusion}\label{sec:conclusion}

We built and evaluated an adaptive exercise recommender for C programming entirely from real learner logs: learner models including DKT and an entropy-based identification tree, a sequential decision problem whose reward combines post-test improvement and solving effort, and DQN and PPO agents. The learner models forecast equally well, yet the agents trained against the DKT simulator obtained large gains that did not survive independently fitted worlds and that the tracker also reported when no learning was possible. We showed that such in-simulator gains are, in expectation, the tracker's accumulated drift---learning the model promises itself---which a coherent model without learning never reports and which a re-seeded DKT did not reproduce. Agents trained with an explicit learning model transferred best within the family of latent worlds, but not to neural worlds. For researchers and developers of adaptive recommenders, the practical message is to accompany any simulator-based evaluation with a negative-control world, the drift decomposition and independently fitted worlds, and to treat in-simulator gains of optimised policies with suspicion until these checks are passed. Replication on other datasets, sensitivity analyses of the decision problem and, ultimately, a controlled study with students are needed to establish whether the robust policies help real learners; our pipeline and application are designed to support such work.
